% Problem dossier: VI.01.P2
\subsection*{Problem VI.01.P2 --- Recovering the equation of a parametrized parabola}
\label{prob:vi01-p2}

\paragraph{Problem.}
Assume \(k\) is infinite and set
\[
S=\{(t,t^2):t\in k\}\subseteq\mathbb A_k^2.
\]
Prove
\[
I(S)=(y-x^2).
\]
Deduce that the smallest algebraic set containing \(S\) is the parabola
\[
V(y-x^2).
\]

\ifdefined\FullProblemDossiers
\paragraph{Why this problem matters.}
A parametrization gives points, while an implicit equation gives a zero locus.  This is the
first nontrivial example where the chapter's \(V\)-\(I\) machine converts one description
into the other.

\paragraph{Useful ideas before starting.}
Treat \(f(x,y)\) as a polynomial in \(y\) with coefficients in \(k[x]\), and divide by the
monic polynomial \(y-x^2\).

\paragraph{Guided hints.}
\textbf{Hint 1.} Certainly \(y-x^2\in I(S)\).

\textbf{Hint 2.} Polynomial division gives
\[
f(x,y)=q(x,y)(y-x^2)+r(x),
\]
with \(r\in k[x]\).

\textbf{Hint 3.} Evaluate at \((t,t^2)\) and use that \(k\) is infinite.

\paragraph{Solution.}
The polynomial \(y-x^2\) vanishes on every \((t,t^2)\), hence
\[
(y-x^2)\subseteq I(S).
\]
Conversely, let \(f\in I(S)\).  Since \(y-x^2\) is monic in \(y\), divide in the ring
\(k[x][y]\):
\[
f(x,y)=q(x,y)(y-x^2)+r(x)
\]
for some \(q\in k[x,y]\) and \(r\in k[x]\).  Evaluating on the parametrized set gives
\[
0=f(t,t^2)=r(t)
\]
for every \(t\in k\).  Because \(k\) is infinite, \(r=0\).  Thus every \(f\in I(S)\) is
divisible by \(y-x^2\), proving
\[
I(S)=(y-x^2).
\]
Therefore
\[
V(I(S))=V(y-x^2),
\]
and by the smallest-algebraic-set property this is the algebraic envelope of \(S\).

\paragraph{What happened mathematically.}
Division eliminated the variable transverse to the graph.  The remainder records the only
possible extra equation along the parameter line, and infinitely many parameter values
force that remainder to vanish identically.

\paragraph{Extensions.}
\begin{enumerate}
\item Replace \(t^2\) by an arbitrary polynomial \(g(t)\) and prove
      \(I\{(t,g(t))\}=(y-g(x))\).
\item Treat the parametrized curve \((t,t^2,t^3)\subseteq\mathbb A^3\).
\item Compare with finite fields, where a nonzero remainder may vanish at every parameter
      value.
\end{enumerate}

\paragraph{Provenance.}
New canonical problem.  It develops an implicitization technique appropriate to the
foundational chapter without consuming the legacy morphism or component problems.
\fi
